\documentclass[../main.tex]{subfiles}

\begin{document}
\subsection{Lema de Goursat (cont.)}
\classdate{2026-04-15}

\begin{lemma}[Lema de Goursat 2.0]\label{lemma:lem-goursat-v2}
    Sea \(U \subseteq \mathbb{C}\) un conjunto abierto, \(R \subseteq U\) un rectángulo
    cerrado, \(z_{0}\in \mathrm{int}(R)\) y \(f \colon U\setminus \{z_{0}\}
    \to \mathbb{C}\) una función analítica, tal que

    \begin{equation*}
        \lim_{z \to z_{0}}f(z)(z - z_{0}) = 0
    \end{equation*}

    entonces

    \begin{equation*}
        \int_{\partial R}^{}f(z)\odif{z} = 0.
    \end{equation*}
\end{lemma}

\begin{proof}
    Por el límite, tenemos que \(\forall \varepsilon > 0\medspace \exists \delta > 0\)
    tal que \(0 < \abs{z - z_{0}} < \delta\),

    \begin{equation*}
        \abs{f(z)(z - z_{0})} < \varepsilon\quad \vee\quad \abs{f(z)} <
        \dfrac{\varepsilon}{\abs{z - z_{0}}}.
    \end{equation*}

    Sea \(R^{\prime}\) un cuadrado de lado \(L\) centrado en \(z_{0}\) tal que \(\forall
    z\in R^{\prime}\; \abs{z - z_{0}} < R^{\prime}\). Prologando los lados de
    \(R^{\prime}\)
    para subdividir \(R\) en nueve rectángulos \(R_{1}, R_{2}, \dots, R_{8},
    R^{\prime}\). Y
    por el \zcref{lemma:lem-goursat-v1}, f es analítica en todos los
    rectángulos, distinto
    \(R^{\prime}\), de modo que

    \begin{equation*}
        \int_{\partial R_{1}}^{}f(z)\odif{z} = \int_{\partial
        R_{2}}^{}f(z)\odif{z} = \cdots =
        \int_{\partial R_{8}}^{}f(z)\odif{z} = 0.
    \end{equation*}

    Y como

    \begin{equation*}
        \int_{\partial R}^{}f(z)\odif{z} = \sum_{i = 1}^{8}\int_{\partial
        R_{i}}^{}f(z)\odif{z} + \int_{\partial R^{\prime}}^{}f(z)\odif{z}
    \end{equation*}

    tal que

    \begin{equation*}
        \int_{\partial R}^{}f(z)\odif{z} = \int_{\partial R^{\prime}}^{}f(z)\odif{z}.
    \end{equation*}

    De modo que
    \begin{equation*}
        \abs`\Bigg{\int_{\partial R}^{}f(z)\odif{z}} = \abs`\Bigg{\int_{\partial
        R^{\prime}}^{}f(z)\odif{z}} \leq \int_{\partial
        R^{\prime}}^{}\abs{f(z)}\abs{\odif{z}}
        \leq \int_{\partial R^{\prime}}^{}\dfrac{\varepsilon}{\abs{z -
        z_{0}}}\abs{\odif{z}}.
    \end{equation*}

    pero \(\forall z\in \partial R^{\prime},\; \abs{z - z_{0}} \geq \tfrac{L}{2}\),

    \begin{align*}
        \int_{\partial R^{\prime}}^{}\dfrac{\varepsilon}{\abs{z -
        z_{0}}}\abs{\odif{z}} &\leq
        \int_{\partial R^{\prime}}^{}\dfrac{2}{L}\varepsilon \abs{\odif{z}},\\
        &= \dfrac{2\varepsilon}{L}\int_{\partial R^{\prime}}^{}\abs{\odif{z}},\\
        &= \dfrac{2\varepsilon}{L}(4L),\\
        &= 8\varepsilon.
    \end{align*}

    Por lo tanto,

    \begin{equation}
        \abs{\int_{\partial R}^{}f(z)\odif{z}} < 8\varepsilon\quad \forall
        \varepsilon > 0
    \end{equation}

    de donde se sigue el resultado

    \begin{equation*}
        \int_{\partial R}^{}f(z)\odif{z} = 0.
    \end{equation*}
\end{proof}

Por inducción tenemos

\begin{corollary}[Goursat 2.1]
    Sea \(U \subseteq \mathbb{C}\) un conjunto abierto, \(R \subseteq U\), \(z_{1},
    \dots, z_{n}\in \mathrm{Int}R\) y \( f \colon U\setminus \{z_{1}, \dots, z_{n}\} \to
    \mathbb{C}\) una función analítica tal que

    \begin{equation*}
        \lim_{z \to z_{i}}f(z)(z - z_{i}) = 0,\quad \forall i\in \{1, \dots, n\}.
    \end{equation*}

    Entonces

    \begin{equation*}
        \int_{\partial R}^{}f(z)\odif{z} = 0.
    \end{equation*}
\end{corollary}

\begin{theorem}[Teorema de Cauchy (local)]\label{theo:cauchy-local}
    Sean \(B \subseteq \mathbb{C}\) una bola abierta y \(f \colon B \to \mathbb{C}\) una
    función analítica. Entonces

    \begin{equation*}
        \int_{\gamma}^{}f(z)\odif{z} = 0
    \end{equation*}

    para cualquier curva cerrada \(\gamma\) suave por pedazos, tal que
    \(\gamma \subseteq
    B\).
\end{theorem}

\begin{proof}
    Usando los primeros dos puntos de la \zcref{prop:equiv-affirmations},
    basta mostrar que
    existe \(F \colon B \to \mathbb{C}\) una función analítica tal que
    \(F^{\prime} = f\).

    Sea \(z_{0}\in  B\) el centro de \(B\) (\(B_{r}(z_{0})\) para algún \(r >
    0\)) y para
    \(z\in B\) definimos las curvas \(\gamma_{1} \colon [a, b] \to
    \mathbb{C}\) dada por la
    unión de un segmento horizontal y un segmento vertical que une a
    \(z_{0}\) con \(z\) y, \(\gamma \colon [c, d] \to
    \mathbb{C}\) la unión de un segmento vertical y uno horizontal que une a
    \(z_{0}\) con
    \(z\), tal que \(-\gamma_{2}\) la curva que hace el recorrido en sentido opuesto.

    De esta manera, por el lema de Goursat,

    \begin{align*}
        0 = \int_{\partial R}^{}f(z)\odif{z} &= \int_{\gamma_{1}}^{}f(z)\odif{z} -
        \int_{\gamma_{2}}^{}f(z)\odif{z},\\
        \int_{\gamma_{1}}^{}f(z)\odif{z} &= \int_{\gamma_{2}}^{}f(z)\odif{z}.
    \end{align*}

    Definimos

    \begin{equation*}
        F(z) = \int_{\gamma_{1}}^{}f(z)\odif{z} = \int_{\gamma_{2}}^{}f(z)\odif{z}.
    \end{equation*}

    Con base en esta definición podemos calcular la derivada de \(F\) respecto a \(x\)
    usando la definición,

    \begin{align*}
        \pdv{F(z)}{x} &= \lim_{h\to 0}\dfrac{F(z + h) - F(z)}{h},\\
        &= \lim_{h\to 0}\dfrac{1}{h}\int_{z}^{z + h}f(z)\odif{z},\\
        &= \lim_{h\to 0}\dfrac{1}{h}\int_{0}^{h}f(z + t)\odif{t} = f(z),
    \end{align*}

    donde para calcular \(F(z)\) consideramos la curva \(\gamma_{1}\)
    mencionada, mientras
    que para calcular \(F(z + h)\) usamos la unión de \(\gamma_{1}\) con un segmento
    horizontal de \(z\) a \(z + h\), donde la resta \(F(z + h) - F(z)\)
    representa entonces
    la integral de \(f\) a lo largo del segmento horizontal que podemos parametrizar por
    \(\gamma_{1}(t) = z + t\). Y, finalmente, como \(f(z)\) es analítica,
    \(\pdv{F(z)}{x}\)
    es continua para toda \(z\).

    De manera análoga (ver \zcref{sec:cr-compleja}),

    \begin{align*}
        \pdv{F(z)}{y} &= \lim_{h\to 0}\dfrac{F(z + ih) - F(z)}{h},\\
        &= \lim_{h\to 0}i \int_{0}^{h}f(z + it)\odif{t},\\
        &= i f(z).
    \end{align*}

    De donde se obtiene que

    \begin{equation*}
        \pdv{F(z)}{x} = i \pdv{F(z)}{y},
    \end{equation*}

    i.e. se cumple cauchy-riemann en su versión compleja. Finalmente, \(F\)
    es analítica y

    \begin{equation*}
        F^{\prime}(z) = \pdv{F(z)}{x} = f(z).
    \end{equation*}
\end{proof}
\end{document}
